% Generated draft: requires mathematical and source review.
\documentclass[11pt]{article}
\usepackage[margin=25mm]{geometry}
\usepackage{fontspec}
\setmainfont{Latin Modern Roman}
\usepackage{amsmath,amssymb,mathtools,unicode-math}
\setmathfont{Latin Modern Math}
\usepackage{enumitem,needspace,xurl,hyperref}
\setlength{\parindent}{0pt}
\setlength{\parskip}{5pt}
\setlength{\emergencystretch}{3em}
\newcommand{\N}{\mathbb{N}}
\newcommand{\Z}{\mathbb{Z}}
\newcommand{\Q}{\mathbb{Q}}
\newcommand{\R}{\mathbb{R}}
\newcommand{\C}{\mathbb{C}}
\newcommand{\F}{\mathbb{F}}
\newcommand{\A}{\mathbb{A}}
\newcommand{\PP}{\mathbb{P}}
\newcommand{\E}{\mathbb{E}}
\newcommand{\eps}{\varepsilon}
\newcommand{\dd}{\,\mathrm{d}}
\begin{document}
\section*{UnsolvedMath ID: 9700028}
\textbf{Source problem code:} catalogue entry
\subsection{Definitions and mathematical statement}
A scale-invariant random spatial network is a random family of spatial routes \(\mathcal R(z,z')\) between planar points, subject to the source's consistency and finiteness axioms. Its central symmetry requirement is joint invariance under Euclidean similarities: for every \(h(x)=cRx+a\), \(c>0\), the finite collections
\[
\bigl(h\mathcal R(z_i,z'_i)\bigr)_i
\quad\text{and}\quad
\bigl(\mathcal R(hz_i,hz'_i)\bigr)_i
\]
have the same law. This is a distributional statement for routes, not a demand that each realized network be fixed by every scaling.

Core finiteness requirements concern finite route length, finite expected length between unit-separated endpoints, and finite expected length intensity in a network sampled at Poisson points. For example, with an independent unit-intensity Poisson set \(\Pi\), the sampled network is
\[
\mathcal S(\Pi)=\bigcup_{z,z'\in\Pi}\mathcal R(z,z').
\]
The exact route-consistency, measurability, and long-route intensity axioms defining SIRSN in the linked paper are not supplied; the displayed core is not substituted for a complete definition.

Open Problem 28 asks to construct a SIRSN that is visually realistic, meaning it does not look very different from a real-world road network. No reference data, statistical distance, tolerance, or formal visual predicate is specified. Hence the task is a model-design programme with qualitative assessment, not a precise extra axiom that can be inferred from scale invariance. Constructing a technically valid SIRSN alone does not certify the requested realism.

The supplied research summary reports Poisson-line constructions and later metric-structure progress, while explicitly not identifying a construction documented to meet Aldous's visual goal. Those technical construction claims and their exact SIRSN hypotheses need their cited papers; no visual judgment or success claim is made here.
\subsection{Short English statement}
Construct a technically valid SIRSN whose appearance also matches the intended qualitative road-network benchmark. The source defines no objective realism test, so known scale-invariant constructions are only partial progress.
\subsection{Sources}
\begin{itemize}
\item Local UnsolvedMath registry, permanent ID \(9700028\), exact source code \texttt{catalogue entry}, ``A visually realistic SIRSN.'' Proposer David Aldous.
\item Aldous, \emph{Open problems}, page \texttt{sirsn.html}, linked paper Open Problem 28; index \url{https://www.stat.berkeley.edu/~aldous/Research/OP/index.html}. Recorded access 2026-07-29; HTML text with linked-primary-document extraction.
\item Supplied PARTIAL-PROGRESS summary names Kendall (2017), Kahn (2016), Kendall (2020), and Blanc--Curien--Kahn (2024/2025) for Poisson-line models and refinements. Complete titles, publication details, and exact axiom scope are not supplied.
\end{itemize}
\end{document}
